\section{Thermomechanical foundations}\label{sec:foundations}
This section collects the balance laws, the thermodynamic structure of linear thermoelasticity, the reasons Fourier's law fails for fast processes, and the heat-flux template that organizes the rest of the review. The presentation is self-contained, so later sections can cite equation numbers instead of rederiving them.

\subsection{Balance laws for small disturbances}
Let a body occupy a region $\Omega$ and let $\bu(\bx,t)$ be the displacement from a reference state at uniform temperature $T_0$ and zero stress. For small displacement gradients the strain is $\varepsilon_{ij}=\tfrac12(u_{i,j}+u_{j,i})$. The local forms of the balance of linear momentum and of energy are
\begin{align}
\rho\,\ddot u_i &= \sigma_{ij,j}+\rho f_i, \label{eq:mom}\\
\rho\,\dot e &= \sigma_{ij}\dot\varepsilon_{ij}-q_{i,i}+\rho r, \label{eq:first}
\end{align}
where $\rho$ is the (reference) density, $f_i$ the body force per unit mass, $e$ the specific internal energy, $q_i$ the heat flux and $r$ the heat supply per unit mass. Balance of angular momentum makes the stress symmetric when no couple stresses act. The second law, in the Clausius--Duhem form used by Coleman and Noll \citep{Coleman1963},
\begin{equation}
\rho\,\dot\eta \;\ge\; \frac{\rho r}{T}-\Bigl(\frac{q_i}{T}\Bigr)_{,i}, \label{eq:clausius}
\end{equation}
with specific entropy $\eta$ and absolute temperature $T$, restricts the constitutive relations; Liu \citep{Liu1972} introduced Lagrange multipliers to impose the entropy principle on the constitutive relations, in place of case-by-case manipulation. Eliminating $r$ between \eqref{eq:first} and \eqref{eq:clausius} and introducing the Helmholtz free energy $\psi=e-T\eta$ gives
\begin{equation}
-\rho\bigl(\dot\psi+\eta\dot T\bigr)+\sigma_{ij}\dot\varepsilon_{ij}-\frac{q_i T_{,i}}{T}\;\ge\;0. \label{eq:CDpsi}
\end{equation}
For a material with $\psi=\psi(\varepsilon_{ij},T)$, \eqref{eq:CDpsi} must hold for arbitrary $\dot\varepsilon_{ij}$ and $\dot T$, so the coefficients of those rates vanish,
\begin{equation}
\sigma_{ij}=\rho\,\frac{\partial\psi}{\partial\varepsilon_{ij}},\qquad \eta=-\frac{\partial\psi}{\partial T},\label{eq:coleman}
\end{equation}
and the remaining term gives the \emph{heat-conduction inequality}
\begin{equation}
q_iT_{,i}\;\le\;0.\label{eq:cond-ineq}
\end{equation}
Fourier's law $q_i=-k_{ij}T_{,j}$ with positive-semidefinite $k_{ij}$ satisfies \eqref{eq:cond-ineq} identically, as does any law whose flux is, to leading order, opposed to the temperature gradient. Here lies the difficulty with lagged and relaxed laws (\cref{sec:beyond}).

\subsection{Linear thermoelasticity}
Expand $\rho\psi$ to second order about $(\varepsilon_{ij},T)=(0,T_0)$, with $\theta=T-T_0$:
\begin{equation}
\rho\psi=\tfrac12C_{ijkl}\varepsilon_{ij}\varepsilon_{kl}-\beta_{ij}\varepsilon_{ij}\theta-\frac{\rho c}{2T_0}\theta^2,\label{eq:psi}
\end{equation}
where $C_{ijkl}$ are the isothermal elastic moduli, $\beta_{ij}$ the thermal-stress tensor and $c$ the specific heat at constant strain. The constant and linear terms in $\theta$ are dropped, since they only fix the reference values of the energy and the entropy. Equations~\eqref{eq:coleman} give
\begin{equation}
\sigma_{ij}=C_{ijkl}\varepsilon_{kl}-\beta_{ij}\theta,\qquad \rho\eta=\beta_{ij}\varepsilon_{ij}+\frac{\rho c}{T_0}\theta.\label{eq:stress}
\end{equation}
The first of these is the Duhamel--Neumann relation. For an isotropic solid $C_{ijkl}=\lambda\delta_{ij}\delta_{kl}+\mu(\delta_{ik}\delta_{jl}+\delta_{il}\delta_{jk})$ and $\beta_{ij}=\beta\delta_{ij}$ with $\beta=(3\lambda+2\mu)\alpha_T$, because free thermal expansion $\varepsilon_{ij}=\alpha_T\theta\,\delta_{ij}$ must produce zero stress: $\sigma_{ij}=(3\lambda+2\mu)\alpha_T\theta\,\delta_{ij}-\beta\theta\,\delta_{ij}=0$. The energy balance $\rho T\dot\eta=-q_{i,i}+\rho r$, which is the entropy form of \eqref{eq:first} once the stress power is accounted for by the free energy, linearized about $T_0$ with the second of \eqref{eq:stress}, gives
\begin{equation}
\rho c\,\dot\theta+T_0\,\beta_{ij}\dot\varepsilon_{ij}=-q_{i,i}+\rho r. \label{eq:en}
\end{equation}
Together with \eqref{eq:mom}, \eqref{eq:stress} and a heat-flux law, \eqref{eq:en} closes the system. Each coefficient has a direct physical meaning. The product $\rho c$ is the heat capacity per unit volume at constant strain; $T_0\beta^2/(\rho c)$ is the difference between the adiabatic and isothermal longitudinal moduli,
\begin{equation}
\lambda_{\rm ad}+2\mu_{\rm ad}=\lambda+2\mu+\frac{T_0\beta^2}{\rho c},
\qquad\text{equivalently}\qquad c_{\rm ad}^2=c_1^2\,(1+\varepsilon),\quad \varepsilon=\frac{\beta^2T_0}{\rho^2c\,c_1^2}. \label{eq:adiabatic}
\end{equation}
\cref{tab:materials} gives the coupling parameter $\varepsilon$ and the characteristic frequency $\omega_*=\rho c c_1^2/k$ for representative solids. The coupling is of order $10^{-2}$ for common metals and far smaller for low-expansion materials. The characteristic frequency lies well above the ultrasonic range for bulk solids ($10^{10}$--$10^{13}$ Hz in \cref{tab:materials}). The thermal frequency that matters falls into the megahertz--gigahertz range only when waves are confined to micro-scale structures, where the Zener frequency $f_Z\propto h^{-2}$ takes over (\cref{sec:applications}).
\begin{table}[ht]\centering\small
\caption{Thermoelastic coupling $\varepsilon$ and characteristic frequency $f_*=\omega_*/2\pi$ computed from rounded handbook-type values (300\,K) of density, specific heat, conductivity, expansion coefficient, Young's modulus and Poisson ratio; the reader should replace them by data for the specimen at hand.}\label{tab:materials}
\input{tables/materials.tex}
\end{table}

\subsection{Beyond Fourier's law}\label{sec:beyond}
\paragraph{Why Fourier's law is parabolic.} Inserting $q_i=-kT_{,i}$ in \eqref{eq:en} with constant coefficients and no coupling gives $\rho c\,\dot\theta=k\nabla^2\theta$. The fundamental solution is a Gaussian, nonzero at every distance for every positive time, so the signal speed is unbounded. Because that Gaussian decays so rapidly, the infinite speed has no observable consequence for slow processes. For processes on the time scale of the carrier relaxation it does.

\paragraph{Relaxation of the flux.} A natural modification, due to Cattaneo and Vernotte \citep{Cattaneo1948,Vernotte1958}, lets the flux relax to its Fourier value with relaxation time $\tau_0$:
\begin{equation}
q_i+\tau_0\dot q_i=-k\,\theta_{,i}.\label{eq:cattaneo}
\end{equation}
With \eqref{eq:en} and no coupling, $k\nabla^2\theta=\rho c(\dot\theta+\tau_0\ddot\theta)$, a damped wave equation with front speed $c_T=\sqrt{k/(\rho c\tau_0)}$. In kinetic terms $\tau_0$ is the mean time between momentum-destroying collisions of the heat carriers, and $c_T$ is of the order of the carrier speed divided by $\sqrt3$ in three dimensions. Phonon hydrodynamics admits a different, ``second sound'', regime in which normal (momentum-conserving) collisions dominate \citep{GuyerKrumhansl1966,GuyerKrumhansl1966II,Chester1963}.

\paragraph{Thermodynamic compatibility via extended irreversible thermodynamics.} The Cattaneo law does not follow from \eqref{eq:clausius} with a free energy that depends only on $(\varepsilon_{ij},T)$: the flux itself must enter the state space. Extended irreversible thermodynamics does that \citep{Jou2010,MuellerRuggeri1998,Mueller1967}, with an entropy that depends also on $q_i$. In one dimension, without deformation, take
\begin{equation}
\eta=\eta_{\rm eq}(e)-\frac{\tau_0}{2\rho kT^2}\,q^2,\qquad \rho\dot e=-q_{,x}
\end{equation}
and an entropy flux $q/T$. Then
\begin{equation}
\rho\dot\eta+\Bigl(\frac qT\Bigr)_{,x}=-\frac{q}{kT^2}\bigl(\tau_0\dot q+kT_{,x}\bigr),
\end{equation}
and for the Cattaneo law the bracket equals $-q$, so the entropy production is $q^2/(kT^2)\ge0$. This is the minimal sense in which the relaxed law is thermodynamically admissible: the entropy carries a quadratic correction in the flux. That correction is of second order in the disturbance, which is why the linear theory does not feel it. In lagged models the corresponding entropy functional must be built model by model. For the time-differential three-phase-lag model, thermodynamic principles impose restrictions on the delay times and, in the high-order case, on the conductivities \citep{Chirita2016,Chirita2017}.

\paragraph{Thermal displacement and the Green--Naghdi route.} Green and Naghdi \citep{Green1991,Green1993} take instead an entropy \emph{balance} with an entropy flux, and postulate constitutive relations in terms of the \emph{thermal displacement} $v$, $\dot v=\theta$ (so $v=\int\theta\,\dd t$). The flux $q_i$ may depend on $\theta_{,i}$ and $v_{,i}$; type II ($q_i=-k^*v_{,i}$) conserves energy, and type III adds Fourier dissipation. In wave problems $v$ is the natural field, since $\hat v=\hat\theta/s$.

\paragraph{Green--Lindsay: rates in the stress and entropy.} Green and Lindsay \citep{Green1972} keep a single temperature but let the free energy and the stress depend on $\dot\theta$, so that the entropy and stress relations carry temperature-rate terms. The original paper expresses this through constitutive coefficients and does not name a relaxation time. The relaxation times $\tau_0,\tau_1$ used in this review are the notation of the later literature for those rate coefficients, and the heat flux is taken in its Fourier form, as is usual for the isotropic case. Solutions of the original theory can show discontinuities at wave fronts. A modification with a strain-rate term has been proposed to remove them \citep{Yu2018ModifiedGL}. \cref{sec:hyperbolic} returns to this theory.

\subsection{The heat-flux template}\label{sec:framework-link}
Linear wave problems with constant coefficients depend on the heat-flux law only through its Laplace-domain relation between flux and temperature gradient. The family is therefore defined
\begin{equation}
q_i+\tau_q\,\dot q_i=-k\,\bigl(\theta_{,i}+\tau_T\,\dot\theta_{,i}\bigr)-k^{*}\bigl(v_{,i}+\tau_v\,\dot v_{,i}\bigr),\qquad \dot v=\theta,\label{eq:template}
\end{equation}
After Laplace transformation with zero initial data
\begin{equation}
\hat q_i=-G(s)\,\hat\theta_{,i},\qquad
G(s)=\frac{k\,(1+\tau_Ts)+k^{*}(1+\tau_vs)/s}{1+\tau_qs}.\label{eq:G}
\end{equation}

\begin{figure}[t!]\centering\small
\input{figs_tikz/fig1_template.tex}
\caption{Original schematic; the setting follows \citep{Biot1956}: the heat-flux template couples the flux to temperature gradients through relaxation and lags. Each theory sets some terms to zero.}\label{fig:template}
\end{figure}

\Cref{fig:template} sketches the template. \Cref{tab:param} lists the parameter choices that reproduce the models expressible in the template. Green--Lindsay theory (\cref{sec:gl}) changes the energy and constitutive equations and so is not a row; for fractional and memory-dependent models $G(s)$ keeps the same structure with $s$ replaced by non-integer powers or by kernel transforms (\cref{sec:memory}).

\begin{table}[ht]\centering\small
\caption{Heat-flux theories as parameter choices in the template \eqref{eq:template}. All entries use the energy balance \eqref{eq:en}.}\label{tab:param}
\begin{tabular}{@{}lccccc l@{}}\toprule
Theory & $k$ & $k^*$ & $\tau_q$ & $\tau_T$ & $\tau_v$ & Source\\\midrule
Classical (Fourier) & $k$ & 0 & 0 & 0 & 0 & \citep{Biot1956}\\
Lord--Shulman (Cattaneo) & $k$ & 0 & $\tau_0$ & 0 & 0 & \citep{Lord1967,Cattaneo1948,Vernotte1958}\\
Dual-phase-lag (DPL), 1st order & $k$ & 0 & $\tau_q$ & $\tau_T$ & 0 & \citep{Tzou1995,Tzou1997}\\
Green--Naghdi II & 0 & $k^*$ & 0 & 0 & 0 & \citep{Green1993}\\
Green--Naghdi III & $k$ & $k^*$ & 0 & 0 & 0 & \citep{Green1991}\\
Three-phase-lag (TPL) & $k$ & $k^*$ & $\tau_q$ & $\tau_T$ & $\tau_v$ & \citep{RoyChoudhuri2007}\\
MGT (flux form) & $k$ & $k^*$ & $\tau_q$ & 0 & 0 & \citep{Quintanilla2019}\\\bottomrule
\end{tabular}
\end{table}

\paragraph{Exact reductions and singular limits.} The function $G(s)$ decides which reductions are regular. The following was verified symbolically (reduction check).
\begin{itemize}
\item \textbf{DPL $\to$ Lord--Shulman:} $\tau_T\to0$ with $\tau_q=\tau_0$ fixed and $k^*=0$. Letting $\tau_q\to0$ with $\tau_T$ fixed gives instead $G=k(1+\tau_Ts)$, a Fourier law with a temperature-rate term (Kelvin--Voigt-like conduction), not Cattaneo.
\item \textbf{TPL $\to$ DPL:} $k^*\to0$. For $k^*\ne0$ the function $G(s)\sim k^*/s$ as $s\to0$, so the low-frequency and long-time response changes qualitatively.
\item \textbf{TPL $\to$ classical:} $k^*\to0$ and all three lags $\to0$.
\item \textbf{GN-II:} $k=0$, all lags zero, $G=k^*/s$; \textbf{GN-III:} all lags zero, $G=k+k^*/s$.
\item \textbf{Large-$s$ behavior.} $G(\infty)=0$ for Lord--Shulman (hyperbolic behavior, finite thermal speed $c_T$); $G(\infty)=k\tau_T/\tau_q$ for DPL, which is diffusive again at high frequency with effective conductivity $k\tau_T/\tau_q$.
\item \textbf{Fractional orders:} the endpoint $\alpha\to2$ of the Caputo family is second order in time and is not Fourier's law (\cref{sec:memory}).
\end{itemize}

\subsection{Energy balance and energy flux}\label{sec:energy}
For the Lord--Shulman model without body forces and heat supply, multiply \eqref{eq:mom} by $\dot u_i$ and \eqref{eq:en} by $\theta/T_0$, and use \eqref{eq:cattaneo}:
\begin{equation}
\frac{\partial\mathcal E}{\partial t}+P_{i,i}=-\frac{q_iq_i}{kT_0},\qquad
\mathcal E=\tfrac12\rho\dot u_i\dot u_i+\tfrac12C_{ijkl}\varepsilon_{ij}\varepsilon_{kl}+\frac{\rho c}{2T_0}\theta^2+\frac{\tau_0}{2kT_0}q_iq_i,\label{eq:energy}
\end{equation}
\begin{equation}
P_i=-\sigma_{ij}\dot u_j+\frac{\theta\,q_i}{T_0}.\label{eq:P}
\end{equation}
Here $\mathcal E$ is positive definite, $P_i$ is the flux associated with it, and the right-hand side is the dissipation. This quadratic (availability-type) energy and its flux $P_i$ are what linear wave theory means by ``energy flux''. They differ from the flux $q_i-\sigma_{ij}\dot u_j$ of the \emph{total} internal energy, whose heat-flux part is first order in the disturbance and averages to zero over a period. For the classical theory the same procedure gives \eqref{eq:energy} without the last term in $\mathcal E$. For other members of the template an energy functional exists under parameter conditions that must be stated model by model: for Green--Naghdi type II the functional is conserved; for the first-order DPL and TPL truncations none is constructed here. Stability of the one-dimensional DPL problem was obtained by a spectral argument \citep{Quintanilla2003DPL}. Thermodynamic compatibility and continuous dependence of the time-differential TPL model are treated in \citep{Chirita2016}. \cref{sec:interfaces} uses \eqref{eq:P} for interface energy partition.

\paragraph{The exact total-energy flux, and the reduced form.}
The availability flux \eqref{eq:P} satisfies a dissipation identity, and it is the flux used for the interface energy partition in \cref{sec:interfaces}. There it is the net flux across the interface, and not the sum of per-mode fractions, that is conserved. It is \emph{not} the flux of
total energy, and the difference matters whenever a fixed control volume is meant rather than a material
element. Add the mechanical power balance to the internal-energy balance, convert the material
derivative with mass conservation, and write $\mathcal E_{\mathrm{tot}}=\rho\bigl(\tfrac12\dot\nu_i\dot u_i+e\bigr)$ for the total energy per unit volume. The balances are read here as spatial balances, with $\rho$, $\sigma_{ij}$ and $q_i$ the current density, Cauchy stress and spatial heat flux. They coincide with the referential quantities of the linearized theory used elsewhere in this review only to first order in the disturbance, and the advection term below is of higher order. The result is the identity
\begin{equation}\label{eq:exactflux}
\frac{\partial \mathcal E_{\mathrm{tot}}}{\partial t}
+\frac{\partial}{\partial x_j}\Bigl(q_j-\sigma_{ji}\dot u_i+\mathcal E_{\mathrm{tot}}\dot u_j\Bigr)
=\rho f_i\dot u_i+\rho r ,
\end{equation}
so that with no body force and no heat supply the outward-positive Eulerian flux is
\begin{equation}\label{eq:exactflux-balance}
F_j=q_j-\sigma_{ji}\dot u_i+\mathcal E_{\mathrm{tot}}\dot u_j,
\qquad
\frac{\partial \mathcal E_{\mathrm{tot}}}{\partial t}+\frac{\partial F_j}{\partial x_j}=0 .
\end{equation}
Three contributions appear: conduction $q_j$, stress power $-\sigma_{ji}\dot u_i$, and advection of
total energy $\mathcal E_{\mathrm{tot}}\dot u_j$. The first two alone make the reduced flux
$J_j=q_j-\sigma_{ji}\dot u_i$ that most of the literature uses, and that first-order small-amplitude
wave calculations use; $J_j$ is the first-order reduction of $F_j$ under the stated small-motion
assumptions, not an exact fixed-volume flux.

None of this displaces \eqref{eq:P}. The two objects answer different questions. $P_i$ is the flux tied
to a positive-definite availability functional whose balance carries the dissipation on the right-hand
side. It is also the flux whose net balance across a welded interface is exact in \cref{sec:interfaces}. $F_j$ is the
conserved flux of total energy in a fixed frame. A comparison of one generalized model's energy transport with
another's means something only when it is stated which of the two is being compared. The
reduced notation common in the literature leaves that open.



\begin{summarybox}[Summary]
The balance laws \eqref{eq:mom}\textendash\eqref{eq:en}, the Coleman--Noll relations \eqref{eq:coleman} and the conduction inequality \eqref{eq:cond-ineq} close the system once a heat-flux law is supplied. Fourier's law is parabolic; relaxed and lagged laws give finite-speed thermal waves. The template \eqref{eq:template} reduces to classical, Lord--Shulman (LS), dual-phase-lag (DPL), three-phase-lag (TPL), GN-II/III and MGT by setting its parameters. The endpoint $\tau_q\to0$ with $\tau_T$ fixed, however, is singular (Kelvin--Voigt-like conduction, not Cattaneo). Among these members, only GN-II has a positive-definite energy functional that is conserved exactly in the stated ideal case; the others dissipate it for nonzero coefficients.
\end{summarybox}
